\chapter{Einleitung}

\section{Ausgangslage und Problemstellung}
In den letzten Jahren hat der Umgang mit Künstliche Intelligenz (KI) zu einer fundamentalen Veränderung in der \ac{IT}-Landschaft geführt \footcite{naveed2023comprehensive}.
Gerade \ac{LLM}-Modelle haben als transformative Technologie die Art und Weise, wie wir arbeiten, wie wir Informationen beschaffen und wie wir Probleme lösen, revolutioniert \footcite{alzubaidi2024researchlandscape}.
Besonders beliebt ist dabei die Integration von bestehenden \ac{LLM}-\acp{API} in eigene Softwarelösungen, da Unternehmen so schnell und kostengünstig von den Vorteilen der \ac{KI} profitieren können \footcite{kucharavy2025impactllmassistants}.
So werden bei Bosch z. B. \ac{LLM}-\acp{API} verwendet, um die \ac{IT}-Infrastrukturen von Anwendungen zu analysieren und zu optimieren.

Das Problem ist dabei, dass die Integration von \ac{LLM}-\acp{API} in Softwarelösungen zwar starke Effizienzsteigerungen ermöglicht, gleichzeitig aber ein oft übersehenes Risiko mit sich bringt. Das Übertragen von kritischen Unternehmensdaten an Dritte korreliert unmittelbar mit unternehmenswirtschaftlichen Erfolgen. Wichtig dabei zu verstehen ist, dass die Verschlüsselung der Anfragen über \ac{TLS} zwar die Integrität und Vertraulichkeit der Daten sichert, nicht aber den Zugriff des Anbieters selbst \footcite{staab2025sokprivacyparadox}.

\section{Zielsetzung und Aufbau der Arbeit}
Um dieses Risiko zu minimieren, gilt es als Aufgabe dieser Arbeit, eine unsichtbare Schicht zwischen der Software und der \ac{API} zu schaffen, mit der kritische Daten maskiert werden. So kann gewährleistet werden, dass die Daten dem \ac{API}-Provider nicht ungeschützt zugänglich sind.

Die Frage, die diese Arbeit dabei beantworten soll, lautet dabei:
\begin{equation}
    \parbox{0.85\textwidth}{\centering \textit{Inwiefern ermöglicht ein hybrider Data-Masking-Ansatz (\ac{Regex} + \ac{NER}) einen effektiven Schutz sensibler \ac{IT}-Infrastrukturdaten bei gleichzeitig hoher Antwortqualität von \ac{LLM}-\acp{API}?}}
\end{equation}
Um dieser nachzugehen, werden zunächst die theoretischen Grundlagen der \ac{LLM}-\ac{API}-Kommunikation und die damit verbundenen Risiken analysiert. Anschließend werden verschiedene Maskierungsansätze untersucht und unter anderem deren Einfluss und Effektivität auf die Antwortqualität untersucht \footcite{zhang2025surveyprivacypreservation}. Dabei schließt die Arbeit mit einer Bewertung der praktischen Einsetzbarkeit unter Berücksichtigung regulatorischer und governance-relevanter Aspekte.